\chapter{Mengalikan Bilangan yang Lebih Besar}\label{ch:g4:mult}

Kelas 3 mengalikan dengan satu angka (\cref{ch:g3:mult}); bab ini
mengalikan dengan bilangan dua angka --- dan menjelaskan ``baris
yang bergeser'' pada cara bersusun yang misterius itu, yang
sebenarnya hanyalah pemecahan menjadi puluhan dan satuan.

\section{Mengalikan dengan puluhan}

\begin{proposition}[Kali 10, 20, 300\,\dots]\label{prop:g4:mult:tens}
Mengalikan dengan $10$ menambahkan satu \omterm{def:g1:counting:zero}{nol}, dengan $100$ dua \omterm{def:g1:counting:zero}{nol}.
Untuk mengalikan dengan $20$, kalikan dengan $2$, lalu dengan $10$:
\[
36 \times 20 = 36 \times 2 \times 10 = 72 \times 10 = 720 .
\]
Begitu pula $14 \times 300 = 14 \times 3 \times 100 = 4\,200$.
\end{proposition}
\admitted

\section{Cara bersusun dengan dua angka}

\begin{example}[Mengapa barisnya bergeser?]\label{ex:g4:mult:why}
Untuk menghitung $23 \times 12$, pecahlah $12$ menjadi $10 + 2$:
\[
23 \times 12 = 23 \times 2 + 23 \times 10 = 46 + 230 = 276 .
\]
Cara bersusun menuliskan tepat kedua \omterm{def:g2:mult:def}{hasil kali} itu, satu di bawah
yang lain --- yang kedua bergeser ke kiri karena ia adalah banyaknya
\emph{puluhan}:
\[
\begin{array}{r}
2\,3 \\
\times\ 1\,2 \\
\hline
4\,6 \\
2\,3\,\cdot \\
\hline
2\,7\,6
\end{array}
\]
(titik itu menandai tempat satuan pada baris yang bergeser, yang
sering dibiarkan kosong).
\end{example}

\begin{omfigure}
\begin{tikzpicture}[scale=0.42]
  \fill[omDef!18] (0,0) rectangle (20,10);
  \fill[omThm!18] (0,10) rectangle (20,12);
  \fill[omDef!34] (20,0) rectangle (23,10);
  \fill[omThm!34] (20,10) rectangle (23,12);
  \draw[thick] (0,0) rectangle (23,12);
  \draw[thick] (20,0) -- (20,12);
  \draw[thick] (0,10) -- (23,10);
  \node[font=\small] at (10,5) {$20 \times 10 = 200$};
  \node[font=\small] at (10,11) {$20 \times 2 = 40$};
  \node[font=\small, rotate=90] at (21.5,5) {$3 \times 10 = 30$};
  \node[font=\small] at (21.5,11) {$6$};
  \node[below, font=\small] at (10,0) {$20$};
  \node[below, font=\small] at (21.5,0) {$3$};
  \node[left, font=\small] at (0,5) {$10$};
  \node[left, font=\small] at (0,11) {$2$};
\end{tikzpicture}

{\small Gambar persegi panjang untuk $23 \times 12$: empat \omterm{def:g2:mult:def}{hasil
kali} yang mudah, $200 + 40 + 30 + 6 = 276$. Cara bersusun
mengelompokkannya menjadi dua baris: $46$ (baris ``$\times 2$'') dan
$230$ (baris ``$\times 10$'').}
\end{omfigure}

\begin{method}[Perkalian bersusun dengan bilangan dua angka]\label{met:g4:mult:column}
Untuk menghitung $457 \times 36$:
\begin{enumerate}
  \item baris pertama: $457 \times 6 = 2\,742$;
  \item baris kedua: $457 \times 3$ \emph{puluhan} $= 1\,371$
        puluhan --- tulislah $1\,371$ bergeser satu tempat ke kiri;
  \item jumlahkan kedua barisnya: $2\,742 + 13\,710 = 16\,452$;
  \item taksirlah untuk memeriksa: $457 \times 36 \approx 500 \times 36
        = 18\,000$? Lebih baik: $450 \times 36$ kira-kira $16\,000$
        --- masuk akal.
\end{enumerate}
\[
\begin{array}{r}
4\,5\,7 \\
\times\ \ \,3\,6 \\
\hline
2\,7\,4\,2 \\
1\,3\,7\,1\,\cdot \\
\hline
1\,6\,4\,5\,2
\end{array}
\]
\end{method}

\begin{example}[Trik menghitung di kepala]\label{ex:g4:mult:tricks}
\begin{itemize}
  \item \emph{Pecahlah dengan cara yang ramah}: $18 \times 11 = 18 \times 10
        + 18 = 198$.
  \item \emph{Pakailah bilangan yang hampir seratus}: $25 \times 99 = 25 \times 100 -
        25 = 2\,475$.
  \item \emph{Lipatduakan lalu bagi dua}: $16 \times 35 = 8 \times 70 =
        560$ (membagi dua satu faktor dan melipatduakan faktor yang
        lain tidak mengubah \omterm{def:g2:mult:def}{hasil kalinya}).
\end{itemize}
\end{example}

\section{Perkalian dalam soal cerita}

\begin{example}\label{ex:g4:mult:problem}
Sebuah sekolah memesan $28$ kotak berisi $145$ lembar kertas.
\begin{enumerate}
  \item Operasinya: $145 \times 28$.
  \item Taksiran: $150 \times 30 = 4\,500$, sedikit terlalu besar
        pada kedua bilangannya.
  \item Bersusun: $145 \times 8 = 1\,160$; $145 \times 2$ puluhan
        $= 2\,900$; seluruhnya $4\,060$.
  \item Kalimat: sekolah itu menerima $4\,060$ lembar kertas.
\end{enumerate}
\end{example}

\section{Latihan}

\begin{exercise}[$\star$]\label{exo:g4:mult:1}
Hitunglah di kepala: $47 \times 10$; \quad $23 \times 100$; \quad
$36 \times 20$; \quad $15 \times 300$; \quad $42 \times 50$.
\end{exercise}

\begin{exercise}[$\star$]\label{exo:g4:mult:2}
Hitunglah $34 \times 21$ dengan memecahnya seperti pada \cref{ex:g4:mult:why}
($34 \times 21 = 34 \times 20 + 34$), lalu periksa secara bersusun.
\end{exercise}

\begin{exercise}[$\star$]\label{exo:g4:mult:3}
Hitunglah secara bersusun: $63 \times 24$; \quad $85 \times 47$; \quad
$126 \times 53$.
\end{exercise}

\begin{exercise}[$\star$]\label{exo:g4:mult:4}
Hitunglah secara bersusun, dengan taksiran lebih dulu: $208 \times 34$;
\quad $517 \times 62$.
\end{exercise}

\begin{exercise}[$\star$]\label{exo:g4:mult:5}
Hitunglah dengan trik dari \cref{ex:g4:mult:tricks}:
$45 \times 11$; \quad $32 \times 99$; \quad $24 \times 45$
(lipatduakan lalu bagi dua --- dua kali kalau kamu mau).
\end{exercise}

\begin{exercise}[$\star$]\label{exo:g4:mult:6}
Sebuah bioskop punya $26$ baris berisi $18$ kursi. Berapa kursi seluruhnya?
\end{exercise}

\begin{exercise}[$\star$]\label{exo:g4:mult:7}
Sebuah truk mengangkut $32$ palet berisi $48$ peti. Berapa peti
seluruhnya? Berikan taksirannya dulu, lalu jawaban yang tepat.
\end{exercise}

\begin{exercise}[$\star$]\label{exo:g4:mult:8}
Carilah kesalahannya: Ani menghitung $57 \times 23$ sebagai
$57 \times 3 = 171$ dan $57 \times 2 = 114$, lalu menjumlahkannya
$171 + 114 = 285$. Jawaban yang benar adalah $1\,311$. Apa yang ia
lupakan?
\end{exercise}

\begin{exercise}[$\star\star$]\label{exo:g4:mult:9}
Seorang tukang kebun menanam $14$ baris berisi $35$ bunga tulip dan
$12$ baris berisi $28$ bunga bakung. Berapa bunga seluruhnya? (Dua \omterm{def:g2:mult:def}{hasil kali}, lalu satu \omterm{def:g1:addition:def}{jumlah}.)
\end{exercise}

\begin{exercise}[$\star\star$]\label{exo:g4:mult:10}
Setiap hari, sebuah toko roti memakai $67$ kg tepung. Kira-kira
berapa tepung dalam satu bulan yang $30$ hari --- taksirlah dengan
bilangan yang \omterm{def:g4:numbers:round}{dibulatkan}, lalu hitunglah dengan tepat.
\end{exercise}

\begin{exercise}[$\star\star$]\label{exo:g4:mult:11}
Jelaskan, dengan gambar persegi panjang seperti pada bab ini,
mengapa $15 \times 12$ dapat dihitung sebagai
$15 \times 10 + 15 \times 2$. Lalu ciptakan pemecahan lain dari
$15 \times 12$ yang juga berhasil.

\end{exercise}
